\documentclass[11pt]{letter}

\usepackage{geometry}
\geometry{margin=1in}
\usepackage{setspace}
\usepackage{hyperref}

\signature{Decory Edwards}
\address{Decory Edwards \\ Johns Hopkins University \\ Department of Economics}

\begin{document}

\begin{letter}{Hiring Committee \\ Cofactor Group \\ Greater Philadelphia Area, PA}

\opening{Dear Hiring Committee,}

I am writing to apply for the Associate Consultant position in Life Sciences at Cofactor Group. I am currently a PhD candidate in Economics at Johns Hopkins University. My training combines quantitative analysis, structured problem solving, and clear communication. I am motivated by problems that require careful framing, disciplined analysis, and actionable recommendations. Cofactor Group’s integrated approach to analytics, market research, and consulting within the life sciences aligns closely with how I approach analytical work.

My research agenda is at the intersection of wealth inequality and macroeconomics. In my job market paper, I build and estimate a structural heterogeneous agent model with returns that differ across households. The core contribution is to show how heterogeneity in returns and income risk jointly shape the wealth distribution and consumption behavior. Relative to closely related work, my framework performs better at matching the lower and middle portions of the wealth distribution. This difference matters quantitatively. Models that focus primarily on returns heterogeneity can fit the top of the wealth distribution closely but tend to generate too much wealth among the bottom half of households. In my framework, income uncertainty generates a precautionary saving motive that produces genuinely low wealth households. This leads to a realistic distribution of marginal propensities to consume and aggregate consumption responses that align with empirical estimates.

A key feature of my research is the use of quantitative analysis to support applied decision making. I identify objectives, collect and structure data from multiple sources, and design analyses that support clear conclusions. I regularly translate complex analysis into written and visual deliverables and work collaboratively to refine recommendations. This work requires precision, curiosity, and the ability to communicate complex ideas clearly, skills that map directly to life sciences consulting engagements.

I am particularly interested in Cofactor Group because of its focus on helping biotechnology and pharmaceutical companies bring innovative medicines to market. The opportunity to work closely with experienced consultants, develop core consulting skills, and gain exposure across the product lifecycle is especially appealing. I value environments that emphasize mentorship, collaboration, and thoughtful analysis.

My economics training has provided a strong foundation in analytical reasoning, quantitative problem solving, and professional communication. I am highly proficient in Excel and comfortable working with the broader Microsoft Office suite to produce client ready analyses and presentations. I work well in team settings and take pride in delivering high quality work under deadlines.

I am also enthusiastic about the prospect of working in a hybrid environment in the Greater Philadelphia area and contributing in person at Cofactor Group’s Wayne office. I see this role as an opportunity to apply my analytical training in a life sciences context while building a long term career in consulting.

I would welcome the opportunity to contribute my skills and perspective to Cofactor Group. Thank you for your time and consideration.

\closing{Sincerely,}

\end{letter}
\end{document}


